\documentclass[10pt,a4paper]{article}
\usepackage[margin=2cm]{geometry}
\usepackage{amsmath,amssymb}
\usepackage{booktabs}
\usepackage{microtype}
\usepackage[hidelinks]{hyperref}
\usepackage{parskip}
\usepackage{xcolor}
\usepackage{titlesec}

\titleformat{\section}{\normalsize\bfseries}{\thesection.}{0.5em}{}
\titleformat{\subsection}{\normalsize\itshape}{\thesubsection.}{0.5em}{}
\titlespacing{\section}{0pt}{8pt}{3pt}
\titlespacing{\subsection}{0pt}{5pt}{2pt}

\definecolor{darkred}{rgb}{0.7,0.0,0.0}
\definecolor{darkgreen}{rgb}{0.0,0.45,0.0}

\title{\textbf{Benchmark Analysis \& Ablation Study}\\[2pt]
  \large Classical vs.\ Quantum vs.\ Classical-Small IQL\\
  + Constant-V + Depth-Matched + Warm-up Ablation
  + Fix A\,+\,B Repair + Fix A Isolation\\
  Hopper-v4 $\cdot$ \texttt{medium-v0} $\cdot$ 100\,k steps}
\author{Yeray Cordero Carrasco \\ QC810 --- Quantum Machine Learning, SDU}
\date{May 2026}

\begin{document}
\maketitle
\thispagestyle{empty}

%──────────────────────────────────────────────────────────────────────────────
\section{Experimental Setup}

Eight variants of Implicit Q-Learning (IQL)~\cite{kostrikov2021} were trained
for 100\,000 gradient steps on \texttt{mujoco/hopper/medium-v0}.
Critic $Q(s,a)$ and actor $\pi(a|s)$ are identical classical MLPs $[256,256]$
across all variants. Only the value network $V(s)$ varies:

\begin{itemize}
  \item \textbf{Classical}: MLP $V(s)$ with hidden dims $[256,256]$,
        \textbf{69\,121} parameters. Seeds $\{0,1,2\}$.
  \item \textbf{Quantum (Warmup)}: DRU circuit $V(s)$, $n_q=8$ qubits,
        $n_L=3$ layers, \textbf{146 parameters}. Layerwise warm-up
        schedule~\cite{skolik2021} activates layers at steps
        $\{0,10\text{k},30\text{k}\}$. Trained with PennyLane
        \texttt{backprop} on \texttt{default.qubit} (GPU).
        Seeds $\{0,1,2\}$.
  \item \textbf{Classical-Small}: single-hidden-layer MLP, width 11,
        \textbf{144 parameters}. Seeds $\{0,1,2\}$.
  \item \textbf{Constant-V}: $V(s)\equiv100$ (frozen scalar),
        \textbf{0 trainable parameters}. The constant matches the affine
        bias $b\approx100$ learned by the quantum circuit across all
        warmup seeds (Section~\ref{sec:scale}). Seeds $\{0,1,2\}$.
  \item \textbf{Classical-Deep}: depth-3 MLP $[8,8,8]$, \textbf{249
        parameters}. Depth-matched to the 3-layer DRU circuit; a fairer
        architectural control than depth-1 Classical-Small.
        Seeds $\{0,1,2\}$.
  \item \textbf{Quantum No-Warmup}: identical DRU circuit (8q, 3L, 146
        params) with all three layers active from step~0.
        Seeds $\{0,\ldots,7\}$ ($n=8$).
  \item \textbf{Quantum Fixed} (\emph{Fix A\,+\,B}): DRU circuit with all
        layers active from step~0 (Fix~B) and affine head initialised at
        $b=374$, $a=55$ (Fix~A), placing $V_0(s)\in[319,429]$ at the
        correct value scale from the first gradient step. $b=374$ is the
        geometric-sum estimate $\bar{r}/(1-\gamma)$ for this dataset;
        $a=55$ matches the standard deviation of the classical $V$.
        Seeds $\{0,1,2\}$.
  \item \textbf{Quantum Fixed-Warmup} (\emph{Fix A only}): DRU circuit
        with the same affine-head initialisation ($b=374$, $a=55$,
        Fix~A) but retaining the original Skolik layerwise warm-up
        schedule. Isolates the effect of correct scale initialisation
        from the effect of warmup removal. Seeds $\{0,1,2\}$.
\end{itemize}

\paragraph{Calibration criteria used throughout.}
We distinguish three increasingly strict criteria for the IQL expectile
calibration $P(A<0)$:
\begin{itemize}
  \item \emph{Direction-correct}: $P(A<0)>0.5$, i.e.\ $V$ sits above the
        median of $Q$.
  \item \emph{On-target}: $P(A<0)\in[0.6,\,0.8]$, within reasonable distance
        of the IQL design target $\tau=0.70$.
  \item \emph{Over-calibrated}: $P(A<0)>0.8$, $V$ above the 80th percentile
        of $Q$.
\end{itemize}
Reported confidence intervals for proportions use the Wilson score interval
at 95\% confidence; tests on returns use Welch's two-sample $t$-test (unequal
variances); tests on proportions use Fisher's exact test (two-sided).

%──────────────────────────────────────────────────────────────────────────────
\section{Results Summary}

\begin{table}[h]
\centering
\caption{Final evaluation return at 100\,k steps. CV = coefficient of
variation. Correct IQL expectile calibration requires
$P(A{<}0)\approx\tau=0.70$. No-warmup uses $n=8$ seeds; all other modes
$n=3$.}
\label{tab:returns}
\setlength{\tabcolsep}{5pt}
\begin{tabular}{lrrrrrr}
\toprule
\textbf{Mode} & \textbf{Mean} & \textbf{Std} & \textbf{CV} &
  \textbf{V-params} & $\mathbb{E}[A]$ & $P(A{<}0)$ \\
\midrule
Classical (69\,121 params)
  & 3505.8 & 16.2  & 0.5\%  & 69\,121 & $-0.94$ & 0.70--0.88 \\
Classical-Deep [8,8,8] (249 params)
  & \textbf{3474.9} & \textbf{8.1} & \textbf{0.2\%} & 249
  & $-1.20$ & $\approx$0.70 \\
\textcolor{darkgreen}{Quantum Fixed (Fix A\,+\,B, 146 params)}
  & \textcolor{darkgreen}{3525.7}
  & \textcolor{darkgreen}{29.3}
  & \textcolor{darkgreen}{0.8\%}
  & \textcolor{darkgreen}{146}
  & \textcolor{darkgreen}{$-0.29$}
  & \textcolor{darkgreen}{0.41--\textbf{0.74}} \\
\textcolor{darkgreen}{Quantum Fixed-Warmup (Fix A only, 146 params)}
  & \textcolor{darkgreen}{3512.1}
  & \textcolor{darkgreen}{41.6}
  & \textcolor{darkgreen}{1.2\%}
  & \textcolor{darkgreen}{146}
  & \textcolor{darkgreen}{$-0.63$}
  & \textcolor{darkgreen}{0.56--0.84} \\
Quantum Warmup (146 params)
  & 3409.1 & 95.2  & 2.8\%  & 146     & $+2.73$ & 0.08--0.42 \\
Quantum No-Warmup (146 params, $n{=}8$)
  & 2794.5 & 817.4 & 29.2\% & 146     & $+4.4$  & 0.03--0.78 \\
\textcolor{darkred}{Constant-V ($V\equiv100$, 0 params)}
  & \textcolor{darkred}{2929.3}
  & \textcolor{darkred}{732.8}
  & \textcolor{darkred}{25.0\%}
  & \textcolor{darkred}{0}
  & \textcolor{darkred}{$+2.88$}
  & \textcolor{darkred}{(inverted)} \\
Classical-Small [11] (144 params)
  & 1957.9 & 1103.6 & 56.4\% & 144    & $-1.89$ & 0.64--0.70 \\
\midrule
\multicolumn{7}{l}{\footnotesize IQL paper reference (1\,M steps):
  $\approx$3\,500 raw return.\quad Constant-V = 83.6\% of Classical.}\\
\bottomrule
\end{tabular}
\end{table}

%──────────────────────────────────────────────────────────────────────────────
\section{Diagnostic Analysis}

\subsection{Quantum V(s) Operates at the Wrong Scale}
\label{sec:scale}

The DRU circuit outputs measurements $\langle Z_0\rangle\in[-1,1]$ scaled by
an affine head $V(s)=a\cdot\langle Z_0\rangle+b$.
Inspection of trained checkpoints reveals:

\begin{center}
\setlength{\tabcolsep}{9pt}
\begin{tabular}{lrrrr}
\toprule
Quantum (Warmup) seed & $a$ & $b$ & $V$ range & $\sigma(V)$ \\
\midrule
0 &  4.96 &  99.8 & $[94.8,\;104.7]$ & 1.74 \\
1 & 64.68 &  94.4 & $[29.7,\;159.1]$ & 19.1 \\
2 & 81.41 &  98.9 & $[17.5,\;180.3]$ & 0.05 \\
\midrule
\multicolumn{5}{l}{\footnotesize Classical V: mean $\approx377$,
  range $[-4,\;434]$, $\sigma\approx55$.}\\
\bottomrule
\end{tabular}
\end{center}

The classical $V(s)$ correctly operates at mean $\approx377$, matching the
geometric-sum approximation $\bar{r}/(1-\gamma)=374$.
All three warmup seeds converge to $b\approx100$, roughly one-quarter of the
correct scale.
This is a structural constraint: the affine head cannot stretch the bounded
measurement output $\langle Z_0\rangle\in[-1,1]$ far enough to reach
$V\approx374$ when initialised at $b\approx0$.

\subsection{Expectile Calibration: Inverted Advantage Sign}

For IQL with $\tau=0.7$, the value function should sit at the 0.7-expectile
of $Q_{\min}(s,a)$, implying $P(A<0)\approx0.70$ over dataset transitions.
Computed directly on 20\,000 dataset samples:

\begin{center}
\setlength{\tabcolsep}{8pt}
\begin{tabular}{lrrrr}
\toprule
Mode / seed & $\bar V$ & $\bar Q$ & $\mathbb{E}[A]$ & $P(A{<}0)$ \\
\midrule
Classical / s0    & 377.0 & 376.4 & $-0.66$ & \textcolor{darkgreen}{0.701} \\
Classical / s1    & 377.2 & 376.5 & $-0.65$ & \textcolor{darkgreen}{0.714} \\
Classical / s2    & 379.3 & 377.7 & $-1.61$ & 0.882 \\
Classical-Deep    & $\sim$377 & $\sim$376 & $\sim{-1.2}$ & $\sim$0.70 \\
\midrule
Quantum Warmup / s0 & 101.2 & 102.3 & $+1.07$ & \textcolor{darkred}{0.288} \\
Quantum Warmup / s1 & 130.4 & 137.8 & $+7.38$ & \textcolor{darkred}{0.420} \\
Quantum Warmup / s2 & 180.3 & 181.7 & $+1.35$ & \textcolor{darkred}{0.076} \\
\bottomrule
\end{tabular}
\end{center}

Classical seeds 0 and 1 land within 0.02 of the IQL target (0.701 and
0.714).
All three warmup seeds have $P(A<0)<0.5$ (95\% Wilson CI for direction-correct
rate: $[0\%,\,56\%]$), with $V$ sitting \emph{below} the mean of $Q$---the
opposite direction from IQL's design.
Seed~2 is particularly revealing: $\sigma(V)=0.05$ (effectively constant,
$V\equiv180$), yet achieves the best return of the original nine runs (3539).
The value network is functioning as a passenger.

\subsection{Constant-V Ablation: Decisive Evidence}
\label{sec:constv}

Replacing $V(s)$ with the frozen scalar $c=100$ (zero trainable parameters,
no backward pass through $V$) yields:

\begin{center}
\setlength{\tabcolsep}{8pt}
\begin{tabular}{lrrr}
\toprule
Constant-V seed & Return & $\mathbb{E}[A]$ & $\mathbb{E}[e^{\beta A}]$ \\
\midrule
0 & 1904.5 & $+2.83$ & 93.9 \\
1 & \textbf{3575.7} & $+3.27$ & 96.6 \\
2 & 3307.6 & $+2.54$ & 93.1 \\
\midrule
Mean $\pm$ Std & $2929.3 \pm 732.8$ & & \\
\bottomrule
\end{tabular}
\end{center}

The constant-V agent achieves 83.6\% of classical performance at zero
parameter cost.
Its advantage statistics ($\mathbb{E}[A]\approx+2.9$,
$\mathbb{E}[e^{\beta A}]\approx94$) are essentially identical to the warmup
quantum agent's ($\mathbb{E}[A]\approx+2.7$, $\mathbb{E}[e^{\beta A}]\approx51$).
Both operate in the same degenerate regime: when $V\equiv c$, the Bellman
target becomes $r+\gamma c$, so $Q\approx r+\gamma c$ and
$A\approx r-c(1-\gamma)$.
The actor then upweights high-reward transitions---a 1-step greedy policy
amounting to reward-weighted behavioural cloning.
The bimodality across constant-V seeds (seed~0 fails at 1904, seeds~1--2
succeed) mirrors the warmup quantum bimodality, confirming variance in this
regime is driven by Q-network initialisation, not by $V$.

\subsection{Depth-Matched Baseline: Classical-Deep [8,8,8]}

\begin{center}
\setlength{\tabcolsep}{8pt}
\begin{tabular}{lrrr}
\toprule
Classical-Deep seed & Return & $\mathbb{E}[A]$ & $P(A{<}0)$ \\
\midrule
0 & 3486.0 & $-1.11$ & $\approx0.70$ \\
1 & 3471.9 & $-1.77$ & $\approx0.70$ \\
2 & 3466.9 & $-0.73$ & $\approx0.70$ \\
\midrule
Mean $\pm$ Std & $3474.9 \pm 8.1$ & & \\
\bottomrule
\end{tabular}
\end{center}

A depth-3 MLP $[8,8,8]$ with 249 parameters achieves 99.1\% of classical
performance with CV\,=\,0.2\%---the most reproducible result across all
non-fixed configurations.
This collapses the initial framing of depth-1 Classical-Small as a fair
parameter-matched control: Classical-Small ($[11]$, 144 params) fails not
because of its parameter budget but because depth-1 MLPs cannot represent
the V-expectile surface on this task.
The DRU circuit has effective depth~3 through its three re-uploading layers;
the correct depth-matched classical baseline succeeds reliably while the
quantum circuit (in its original configuration) does not.

%──────────────────────────────────────────────────────────────────────────────
\section{Warm-up Schedule Ablation: Quantum No-Warmup}

\subsection{Per-Seed Results, $n=8$}

\begin{center}
\setlength{\tabcolsep}{5pt}
\begin{tabular}{clrrrrrr}
\toprule
\textbf{Seed} & \textbf{Config} & \textbf{Return} & $\mathbb{E}[A]$
  & $P(A{<}0)$ & $\|\nabla\theta\|_\text{early}$
  & $\|\nabla\theta\|_\text{late}$ & $\|\nabla\theta\|_\text{max}$ \\
\midrule
0 & Warmup    & 3374.7 & $+1.33$ & 0.288
  & 0.395 &  1.40 &   4.2 \\
0 & No-warmup & 3238.0 & $+31.2$ & \textcolor{darkred}{0.030}
  & 3.158 & 97.37 & 448.4 \\
\midrule
1 & Warmup    & 3313.6 & $+5.60$ & 0.420
  & 0.00004 & 52.25 & 199.2 \\
1 & No-warmup & \textbf{3606.1} & $\mathbf{-3.03}$
  & \textcolor{darkgreen}{\textbf{0.784}}
  & 5.581 & 13.98 &  37.7 \\
\midrule
2 & Warmup    & 3539.0 & $+1.27$ & 0.076
  & 0.551 &  1.83 &  22.8 \\
2 & No-warmup & 3535.2 & $+3.68$ & \textcolor{darkred}{0.181}
  & 6.206 &  9.18 &  14.1 \\
\midrule
3 & No-warmup & 1996.8 & $+3.76$ & \textcolor{darkred}{0.422}
  & 2.788 & 34.26 &  98.0 \\
4 & No-warmup & 2595.5 & $+2.30$ & \textcolor{darkred}{0.410}
  & 2.750 & 28.80 &  62.0 \\
\midrule
5 & No-warmup & 2612.8 & $+0.61$ & 0.575
  & --- & --- &  11.6 \\
6 & No-warmup & 3589.2 & $+4.81$ & \textcolor{darkred}{0.176}
  & --- & --- &  51.4 \\
7 & No-warmup & 1182.1 & $+7.36$ & \textcolor{darkred}{0.394}
  & --- & --- & 187.4 \\
\bottomrule
\end{tabular}
\end{center}
\vspace{3pt}
``Early'' = mean $\|\nabla\theta\|$ at steps 0--10k; ``late'' = steps
30k--100k. Seeds 5--7 were analysed post-hoc without early/late gradient
breakdowns.

\subsection{Statistical Summary, Warmup vs No-Warmup}

\begin{center}
\setlength{\tabcolsep}{5pt}
\begin{tabular}{lrrrrr}
\toprule
\textbf{Config} & \textbf{Mean return} & \textbf{CV}
  & $N(P(A{<}0)>0.5)$ & $N(P{\in}[0.6,0.8])$
  & $N(\|\nabla\theta\|>50)$ \\
\midrule
Warmup ($n=3$)    & $3409.1 \pm  95.2$ &  2.8\% & 0/3 & 0/3 & 1/3 \\
No-Warmup ($n=8$) & $2794.5 \pm 817.4$ & 29.2\% & 2/8 & 1/8 & 5/8 \\
\bottomrule
\end{tabular}
\end{center}

\paragraph{Return.} Welch's $t$-test on the 11 returns gives $t=1.94$,
$p=0.090$ (Cohen's $d=0.80$): the no-warmup mean is substantially below the
warmup mean, but the comparison is borderline rather than conventionally
significant due to the n=3 warmup arm.
A more powerful test compares no-warmup to Constant-V, since the diagnostic
prediction is that without scale correction the quantum agent collapses to the
constant-V regime. Welch's $t$-test on these gives $t=-0.22$, $p=0.84$:
\textbf{no-warmup return is statistically indistinguishable from a frozen
scalar value function.} This is the crispest empirical demonstration that
warmup removal alone does not rescue the quantum value learner.

\paragraph{Calibration.}
Direction-correct rate: 0/3 (warmup) vs 2/8 (no-warmup); on-target rate:
0/3 vs 1/8.
Wilson 95\% CIs are $[0\%,56\%]$ for warmup direction-correct,
$[7\%,59\%]$ for no-warmup direction-correct, and $[2\%,47\%]$ for no-warmup
on-target.
Fisher exact tests do not reject independence at any of these contrasts
($p\geq0.49$): n=3 versus n=8 is too underpowered to detect the difference
formally, but the no-warmup configuration is the only one of the two in which
\emph{any} seed exits the inverted-advantage regime.

\paragraph{Gradient explosion.}
Adopting a fixed threshold $\|\nabla\theta\|_\text{max}>50$ for explosion:
warmup 1/3, no-warmup 5/8. Fisher exact $p=0.55$.
At a more permissive threshold $>10$, all 8 no-warmup seeds qualify against
2/3 warmup seeds (Fisher $p=0.27$). Neither comparison is significant, but
the no-warmup gradient distribution is qualitatively heavier-tailed (max
values 14--448 versus 4--199) and the explosions correlate with calibration
failure: among the 5 no-warmup seeds with $\|\nabla\theta\|>50$, all 5 fail
the on-target criterion.

\subsection{No-Warmup as a Stability/Calibration Trade-off}

The two configurations represent different points on a stability-versus-
calibration trade-off rather than one being strictly preferable.

Warmup suppresses early-phase gradients (0.4--0.6) and yields low return CV
(2.8\%). The cost is that direction-correct calibration is never achieved
(0/3) and on-target calibration is never achieved (0/3): the Skolik schedule
prevents the V-Q bootstrapping loop from finding the trajectory that places
$V$ above $\bar Q$.

No-warmup allows the full circuit to explore the optimisation landscape from
step~0. This raises the gradient explosion rate (5/8 at $>50$, 8/8 at $>10$)
and inflates return CV to 29.2\%, but creates the conditions under which a
minority of seeds reach the direction-correct regime (2/8, 95\% CI
$[7\%,59\%]$).

Neither configuration is reliable in isolation. Their juxtaposition motivates
a different intervention---one that preserves multi-layer gradient exploration
but constrains the optimisation initial condition rather than its dynamics.
This is what Fix~A does (Section~\ref{sec:fixab}).

\subsection{No-Warmup Seed 1: The Only On-Target Default-Init Quantum Run}

Across all 11 default-init quantum seeds (warmup $n=3$ + no-warmup $n=8$),
seed~1 (no-warmup) is the sole run to satisfy the on-target criterion
$P(A{<}0)\in[0.6,0.8]$.
Its affine head converges to $a=32.1$, $b=61.0$, giving $V\in[42.4,\,80.1]$
with $P(A{<}0)=0.784$ and $\mathbb{E}[A]=-3.3$.
Value loss decreases monotonically from 20 to 11, gradient norms are stable
throughout ($5.6\to14\to14$, max 37.7), and $\sigma(V)$ settles at 8.97,
indicating meaningful state-dependent discrimination.

Note that this calibration is correct in expectile \emph{direction}
($V$ above $\bar Q$) but at the wrong absolute scale: $\bar V\approx62$
versus the dataset-correct value $\approx374$.
A well-calibrated policy still results ($\text{return}=3606$), consistent
with Hopper-medium's properties: the actor succeeds whenever advantage signs
are correct, regardless of absolute V magnitude.
This seed demonstrates the circuit \emph{can} find the correct optimisation
trajectory from random initialisation, with empirical probability 1/8
(Wilson CI $[2\%,47\%]$).

\subsection{Why the Warm-up Schedule Is Not the Right Mitigation Here}

The Skolik schedule was validated on variational quantum eigensolvers (VQE)
and supervised quantum models with stationary loss landscapes.
In IQL, the V-loss landscape shifts continuously as Q bootstraps from
$\bar{V}$: activating a new layer abruptly changes $V(s')$, which changes the
Bellman target $r+\gamma\bar{V}(s')$, which in turn changes the Q-loss
landscape.
The phase-activation transitions are therefore not curriculum steps in a fixed
landscape but discontinuous perturbations of a coupled V-Q optimisation.

The empirical evidence is consistent with this account. Warmup seed~1 collapses
to a barren plateau in its restricted 1-layer subspace
($\|\nabla\theta\|_{L=1}\approx4\times10^{-5}$), then receives a destabilising
perturbation when layers~2 and~3 are activated, ultimately reaching
$\|\nabla\theta\|_\text{max}=199$.
No-warmup seed~1---the same initialisation without the schedule---avoids both
failure modes and reaches direction-correct calibration.
The schedule trades one failure mode (multi-layer chaos) for another
(single-layer barren plateau plus target-shift shock at activation
boundaries) rather than solving the underlying instability.

%──────────────────────────────────────────────────────────────────────────────
\section{Fix A\,+\,B Repair and Fix A Isolation}
\label{sec:fixab}

\subsection{Per-Seed Results}

\begin{center}
\setlength{\tabcolsep}{4.5pt}
\begin{tabular}{clrrrrrrr}
\toprule
\textbf{Seed} & \textbf{Config} & \textbf{Return} & $\mathbb{E}[A]$
  & $P(A{<}0)$ & $\mathcal{L}(V)_\text{final}$
  & $\|\nabla\theta\|_\text{max}$ & $b_\text{final}$ & $a_\text{final}$ \\
\midrule
0 & Warmup       & 3374.7 & $+1.07$ & 0.288 &  99.1 &   4.2 &  99.8 &   5.0 \\
0 & No-warmup    & 3238.0 & $+31.1$ & 0.030 & 1130  & 320   &  90.8 &  65.7 \\
0 & \textbf{Fixed}        & 3565.1 & $+5.91$ & 0.412 & 407 & 656 & 455.8 & 123.4 \\
0 & Fixed-Warmup & 3466.1 & $-0.70$ & \textcolor{darkgreen}{0.844} & 44.4 & 109.6 & 401.3 & 3.9 \\
\midrule
1 & Warmup       & 3313.6 & $+5.60$ & 0.420 &  99.1 & 199   &  94.4 &  64.7 \\
1 & No-warmup    & 3606.1 & $-3.03$ & \textcolor{darkgreen}{0.784} &  11.4 &  37.7 &  61.0 &  32.1 \\
1 & \textbf{Fixed}        & 3517.0 & $-0.87$ & \textcolor{darkgreen}{\textbf{0.744}} & \textbf{2.0} & 851 & 378.2 & 58.5 \\
1 & Fixed-Warmup & 3503.4 & $-0.77$ & \textcolor{darkgreen}{0.814} & 1.5 & 326.9 & 383.5 & 6.4 \\
\midrule
2 & Warmup       & 3539.0 & $+1.27$ & 0.076 &  99.1 &  22.8 &  98.9 &  81.4 \\
2 & No-warmup    & 3535.2 & $+3.68$ & 0.181 &  27.7 &  14.1 &  98.0 &  11.3 \\
2 & \textbf{Fixed}        & 3495.0 & $-0.58$ & \textcolor{darkgreen}{0.534} & 17.1 & 441 & 380.6 &   2.6 \\
2 & Fixed-Warmup & 3566.8 & $-0.42$ & \textcolor{darkgreen}{0.555} & 24.0 & 189.4 & 419.7 & 3.5 \\
\bottomrule
\end{tabular}
\end{center}

\subsection{Aggregate Calibration Counts and Return Comparisons}

\begin{center}
\setlength{\tabcolsep}{5pt}
\begin{tabular}{lrrrrr}
\toprule
\textbf{Config} & \textbf{Mean return} & \textbf{CV}
  & $N(P{>}0.5)$ & $N(P{\in}[0.6,0.8])$ & $N(P{>}0.8)$ \\
\midrule
Warmup ($n=3$)              & $3409.1 \pm  95.2$ &  2.8\% & 0/3 & 0/3 & 0/3 \\
No-Warmup ($n=8$)           & $2794.5 \pm 817.4$ & 29.2\% & 2/8 & 1/8 & 0/8 \\
\textcolor{darkgreen}{\textbf{Fixed (A\,+\,B, $n=3$)}}
  & \textcolor{darkgreen}{\textbf{3525.7 $\pm$ 29.3}}
  & \textcolor{darkgreen}{\textbf{0.8\%}}
  & \textcolor{darkgreen}{\textbf{2/3}}
  & \textcolor{darkgreen}{\textbf{1/3}} & 0/3 \\
\textcolor{darkgreen}{Fixed-Warmup (A only, $n=3$)}
  & \textcolor{darkgreen}{$3512.1 \pm 41.6$}
  & \textcolor{darkgreen}{1.2\%}
  & \textcolor{darkgreen}{3/3} & 0/3 & 2/3 \\
\midrule
\textbf{Fix-A pool (Fixed + Fixed-Warmup, $n=6$)}
  & $\mathbf{3519.1 \pm 39.6}$ & 1.1\% & \textbf{5/6} & 1/6 & 2/6 \\
\textbf{Default pool (Warmup + No-Warmup, $n=11$)}
  & $\mathbf{2962.3 \pm 786.7}$ & 26.6\% & \textbf{2/11} & 1/11 & 0/11 \\
\bottomrule
\end{tabular}
\end{center}

\paragraph{Pooled Fix-A vs default-init: the headline statistical result.}
Pooling the two Fix-A configurations against the two default-init
configurations yields the only conventionally significant tests in the study:

\begin{itemize}
  \item \emph{Final return}: Fix-A pool $3519\pm40$ (n=6) vs default pool
        $2962\pm787$ (n=11). Welch's $t=2.34$, $\mathbf{p=0.041}$, Cohen's
        $d=0.87$ (large effect).
  \item \emph{Direction-correct calibration}: Fix-A 5/6 (Wilson CI
        $[44\%,97\%]$) vs default 2/11 (Wilson CI $[5\%,48\%]$). Fisher
        exact two-sided $\mathbf{p=0.034}$.
\end{itemize}

These two tests provide the principal statistical evidence that correct
scale initialisation (Fix~A) is the dominant and effective intervention.

\paragraph{Within-cell tests are underpowered.}
Pairwise Welch tests on returns at n=3 versus n=3 (or n=3 versus n=8) do not
reach significance even when point estimates differ substantially:

\begin{center}
\setlength{\tabcolsep}{8pt}
\begin{tabular}{lrr}
\toprule
Comparison & Welch $p$ & Cohen's $d$ \\
\midrule
Fixed vs Warmup                     & 0.22  & 1.35 \\
Fixed vs No-Warmup                  & 0.050 & 0.95 \\
Fixed vs Fixed-Warmup               & 0.73  & 0.31 \\
Fixed-Warmup vs Warmup              & 0.26  & 1.15 \\
Fixed-Warmup vs No-Warmup           & 0.053 & 0.93 \\
Warmup vs No-Warmup                 & 0.090 & 0.80 \\
No-Warmup vs Constant-V             & 0.84  & $-0.15$ \\
\bottomrule
\end{tabular}
\end{center}

The large effect sizes confirm that the absent significance is a power
problem, not a null result. The pooled Fix-A vs default test is what allows
the diagnostic claim to be made statistically.

\subsection{Fix-A\,+\,B (Fixed) Seed 1: First Classical-Quality Quantum
            Value Function}

Quantum-Fixed seed~1 satisfies all four classical-quality criteria
simultaneously:
\begin{itemize}
  \item $P(A{<}0) = 0.744$---above the IQL target $\tau=0.70$ and within
        0.04 of the best classical seeds (0.701, 0.714).
  \item $\mathcal{L}(V)_\text{final} = 2.0$---within a factor of $1.8\times$
        of classical's $1.12\pm0.34$.
  \item $\mathbb{E}[A] = -0.87$---negative, correct direction, magnitude
        consistent with classical ($-0.64$ to $-1.61$).
  \item $b_\text{final}=378.2$, $a_\text{final}=58.5$, giving
        $V\in[370.8,\,384.9]$ centred at $\bar V=378$---matching the
        classical mean of 377 and the correct-scale criterion, with
        state-dependent slope within 1.1$\times$ of classical's
        $\sigma\approx55$.
\end{itemize}

This is the only seed across the 17 quantum runs in this study (3+8+3+3) to
satisfy all four criteria simultaneously.
The next-closest is Fixed-Warmup seed~1, which matches on $\mathcal{L}(V)$
and $b$ but is over-target on $P(A{<}0)$ and has $a=6.4$ (far smaller than
classical's slope), indicating the circuit's state-dependence is suppressed
when warmup is retained.
This proof-of-concept establishes that the DRU architecture is not inherently
incompatible with IQL value learning; the failures were in the training
configuration.

\subsection{Failure Mode A: Initial Value-Loss Shock (Fixed Seed 0)}

With $V_0(s)\approx374$ and $Q_0(s,a)\approx0$ (random critic initialisation),
the expectile loss at step~0 is approximately
$\mathcal{L}(V)_0 \approx \tau\cdot374^2\approx98\,000$.
Gradient clipping at unit max-norm preserves gradient direction but not
magnitude, so the optimiser takes approximately 10\,000 unit-norm steps all
pointing in the same direction before $Q$ has bootstrapped to the correct
scale. The result is systematic overshoot: $b_\text{final}=455$, 22\% above
the target 374.

This is a third failure mode, distinct from those previously identified:
(i)~warmup's 1-layer barren plateau; (ii)~no-warmup's multi-layer chaotic
initialisation; (iii)~\textbf{cold-start value-loss shock} from the V-Q
bootstrap gap.
The natural remedy is \textbf{Fix~C}: freeze the value gradient for the first
1\,000--2\,000 steps, allowing $Q$ to bootstrap toward
$r+\gamma b_0\approx374$ before $V$ begins updating. This would reduce
$\mathcal{L}(V)_0$ from $\sim$98\,000 to $\sim$0, eliminating the overshoot
without any architectural change.
Fix~C is untested within the scope of this report.

\subsection{Failure Mode B: Correct Scale, Absent State-Dependence
            (Fixed Seed 2)}

Seed~2 correctly learns the scale ($b_\text{final}=380.6\approx374$) but
converges to near-zero slope ($a_\text{final}=2.6$), producing
$V\in[378.1,\,383.2]$---a constant at the right level.
$P(A{<}0)=0.534$: direction-correct (above 0.5) but below the on-target band.
Gradient norms collapse to near zero late in training
($\|\nabla\theta\|_\text{late}=2.1$), indicating convergence to a flat local
minimum in circuit parameter space.
The circuit has solved the scale problem but not the state-discrimination
problem.
Fix~C combined with a slightly larger initial $a$ may improve this seed by
providing a better-conditioned starting point for the slope parameter.

\subsection{Fix A Isolation: What Does Removing the Warmup Add?}

Comparing the two Fix-A variants directly:

\begin{itemize}
  \item Returns are statistically indistinguishable: Fixed $3526\pm29$
        vs Fixed-Warmup $3512\pm42$, Welch $p=0.73$, $d=0.31$.
  \item \textbf{Advantage sign}: all three Fixed-Warmup seeds achieve
        $\mathbb{E}[A]<0$ ($-0.70$, $-0.77$, $-0.42$), compared to 2/3 for
        Fixed (seeds 1 and 2 negative, seed 0 positive at $+5.91$).
        Both Fix-A variants outperform default-init in this criterion
        (warmup: 0/3 negative; no-warmup: 2/8 negative).
  \item \textbf{Magnitude}: Fixed-Warmup advantage means cluster tightly
        ($\mathbb{E}[A]\in[-0.42,\,-0.77]$), while Fixed shows larger spread
        ($[-0.58,\,+5.91]$). The warmup schedule's phase-by-phase activation
        appears to suppress the initial V-Q shock that causes seed~0's
        large positive $\mathbb{E}[A]$ in Fixed.
  \item Direction-correct calibration ($P(A{<}0)>0.5$): Fixed 2/3 vs
        Fixed-Warmup 3/3 (Fisher exact $p=1.0$, indistinguishable at this
        sample size).
  \item Over-calibration ($P>0.8$): Fixed 0/3 vs Fixed-Warmup 2/3.
        Fisher exact $p=0.40$---suggestive but not significant.
\end{itemize}

The negative $\mathbb{E}[A]$ in all Fixed-Warmup seeds confirms that Fix~A
(correct scale initialisation) achieves the correct IQL optimisation
direction reliably, regardless of whether the warmup schedule is retained.
The descriptive pattern is consistent with a hypothesis that retaining the
warmup schedule alongside Fix~A causes the multi-layer circuit to begin
optimising with $b$ already near 374 and $Q$ still bootstrapping, allowing
$V$ to overshoot $Q$ before $Q$ rises---producing the observed
over-calibration in 2/3 Fixed-Warmup seeds and the tight negative $\mathbb{E}[A]$
cluster. However, with 3 seeds per arm,
this hypothesis is not statistically testable; it is a candidate for the
follow-up seed expansion to $n=8$.

A more conservative reading of the current data: \textbf{Fix~A is the
intervention that does the work}. Removing the warmup additionally (Fix~B)
produces a small, statistically undetectable improvement in calibration
centring and a small, statistically undetectable degradation in gradient
peaks (Fixed seeds reach $\|\nabla\theta\|_\text{max}$ up to 851; Fixed-Warmup
peaks at 327). Both variants explode at the cold-start step and recover.
Choosing between Fixed and Fixed-Warmup as the recommended configuration
requires more seeds.

\paragraph{Implication for the no-warmup analysis.}
The fact that Fix-A combined with the warmup schedule still produces 5/6
direction-correct calibration across the two Fix-A variants reframes the
no-warmup result. The warmup schedule on its own---without correct scale
initialisation---is incompatible with IQL value learning (Section 4).
But once the affine bias is correctly initialised, the warmup schedule is
no longer fatal; it is merely suboptimal. The dominant problem in the
original report was scale, not schedule.

%──────────────────────────────────────────────────────────────────────────────
\section{Statistical Significance Summary}
\label{sec:stats}

The conclusions of this study rest on the following statistical evidence,
ordered by strength.

\paragraph{Conventionally significant ($p<0.05$).}
\begin{enumerate}
  \item \textbf{Fix-A pool versus default-init pool, return.}
        Welch's $t=2.34$, $\mathbf{p=0.041}$, $d=0.87$, $n_1{=}6$,
        $n_2{=}11$. Mean returns $3519\pm40$ versus $2962\pm787$.
  \item \textbf{Fix-A pool versus default-init pool, direction-correct
        calibration.} Fisher exact two-sided $\mathbf{p=0.034}$,
        5/6 (Wilson CI $[44\%,97\%]$) versus 2/11 (Wilson CI
        $[5\%,48\%]$).
\end{enumerate}

\paragraph{Strong negative result.}
\begin{enumerate}
  \setcounter{enumi}{2}
  \item \textbf{No-Warmup return is statistically indistinguishable from
        Constant-V.} Welch's $t=-0.22$, $p=0.84$, $d=-0.15$. This
        confirms that warmup removal alone produces a quantum agent whose
        return distribution matches a frozen scalar baseline, supporting the
        diagnostic claim that scale correction (not schedule modification)
        is the active intervention.
\end{enumerate}

\paragraph{Borderline ($0.05<p<0.10$, large effect size).}
\begin{enumerate}
  \setcounter{enumi}{3}
  \item \emph{Fixed vs No-Warmup return}: $p=0.050$, $d=0.95$.
  \item \emph{Fixed-Warmup vs No-Warmup return}: $p=0.053$, $d=0.93$.
  \item \emph{Warmup vs No-Warmup return}: $p=0.090$, $d=0.80$.
        Treated descriptively rather than significantly throughout.
\end{enumerate}

\paragraph{Underpowered (large effect, no significance).}
Pairwise tests at n=3 versus n=3 (Fixed vs Warmup, Fixed-Warmup vs Warmup,
Fixed vs Fixed-Warmup, gradient explosion comparisons, on-target calibration
comparisons) do not reach significance despite Cohen's $d$ values up to 1.35.
These should be read as observations consistent with the hypothesised effects,
not as confirmed effects.

\paragraph{Wilson 95\% confidence intervals for key proportions.}

\begin{center}
\setlength{\tabcolsep}{8pt}
\begin{tabular}{lll}
\toprule
Quantity & Point estimate & 95\% Wilson CI \\
\midrule
Warmup direction-correct                & 0/3 = 0\%    & $[0\%,\,56\%]$ \\
No-Warmup direction-correct             & 2/8 = 25\%   & $[7\%,\,59\%]$ \\
No-Warmup on-target $[0.6,0.8]$         & 1/8 = 12.5\% & $[2\%,\,47\%]$ \\
Fixed direction-correct                 & 2/3 = 67\%   & $[21\%,\,94\%]$ \\
Fixed on-target                         & 1/3 = 33\%   & $[6\%,\,79\%]$ \\
Fixed-Warmup direction-correct          & 3/3 = 100\%  & $[44\%,\,100\%]$ \\
Fixed-Warmup over-calibrated $({>}0.8)$ & 2/3 = 67\%   & $[21\%,\,94\%]$ \\
\textbf{Fix-A pool direction-correct}   & \textbf{5/6 = 83\%} & $\mathbf{[44\%,\,97\%]}$ \\
Default-init direction-correct          & 2/11 = 18\%  & $[5\%,\,48\%]$ \\
\bottomrule
\end{tabular}
\end{center}

The CIs make clear that single-cell rates at n=3 are not pinned down within
useful bounds, and the strongest claim the data supports is the contrast
between the pooled Fix-A and default-init proportions.

%──────────────────────────────────────────────────────────────────────────────
\section{Conclusions}

\textbf{(1) The quantum $V(s)$ does not function as an IQL value function in
its default configuration.}
All three warmup seeds produce $P(A{<}0)\in[0.08,\,0.42]$ against the required
$\approx0.70$, with $V$ sitting below the mean of $Q$ rather than above it.
The structural cause is an affine output head that cannot reach the correct
value scale ($\approx374$) from a zero-centred initialisation.

\textbf{(2) The default quantum system degenerates to Q-only reward-weighted
behavioural cloning.}
Constant-V ($V\equiv100$, zero parameters) achieves 83.6\% of classical
performance and produces advantage statistics indistinguishable from the
warmup quantum agent. Both operate in the regime $Q\approx r+\gamma c$,
$A\approx r-c(1-\gamma)$, with the actor upweighting high-reward transitions.
The crispest demonstration is the Welch test on No-Warmup versus Constant-V
returns ($p=0.84$): removing the warmup without correcting scale produces a
quantum agent indistinguishable from a frozen scalar.

\textbf{(3) Competitive return in the default configuration is a property of
the task, not the circuit.}
Hopper-medium's dense reward, positive correlation with locomotion quality,
and narrow dataset distribution make reward-weighted BC competitive with full
IQL.

\textbf{(4) Depth is the relevant inductive bias, not parameter count.}
Classical-Deep $[8,8,8]$ (249 params, depth~3) achieves 99.1\% of classical
performance reliably (CV\,=\,0.2\%); Classical-Small $[11]$ (144 params,
depth~1) fails with CV\,=\,56.4\%.
The DRU circuit has effective depth~3; the fair classical comparison is
depth-matched, and depth-3 classical wins decisively.

\textbf{(5) Warmup removal alone does not rescue the quantum agent.}
At $n=8$, no-warmup produces mean return $2794\pm817$ (CV\,=\,29.2\%), 12\%
below warmup ($p=0.090$, $d=0.80$) and \emph{statistically indistinguishable
from Constant-V} ($p=0.84$).
Direction-correct calibration is achieved in 2/8 seeds (Wilson CI
$[7\%,59\%]$), on-target in 1/8 (CI $[2\%,47\%]$).
At threshold $\|\nabla\theta\|_\text{max}>50$, gradient explosion occurs in
5/8 no-warmup runs versus 1/3 warmup runs (Fisher $p=0.55$, descriptive).
The warmup schedule was designed for stationary loss landscapes (VQE,
supervised quantum models) and is incompatible with IQL's continuously
shifting V-Q bootstrapping loop, but removing it without scale correction
makes things worse, not better.

\textbf{(6) Fix~A (correct scale initialisation) is the dominant intervention,
with conventionally significant statistical support.}
Pooling the two Fix-A configurations (Fixed and Fixed-Warmup, $n=6$) against
the two default-init configurations (Warmup and No-Warmup, $n=11$):

\begin{itemize}
  \item Return: $3519\pm40$ vs $2962\pm787$, Welch $\mathbf{p=0.041}$,
        $d=0.87$.
  \item Direction-correct calibration: 5/6 vs 2/11, Fisher exact
        $\mathbf{p=0.034}$.
\end{itemize}

These are the only conventionally significant comparisons in the study.
They establish that the DRU architecture is not inherently incompatible with
IQL value learning---the failures were in the training configuration.

\textbf{(7) Fix~B (warmup removal in addition to Fix~A) is statistically
undetectable at this sample size.}
Within the Fix-A pool, Fixed and Fixed-Warmup are statistically
indistinguishable in return (Welch $p=0.73$, $d=0.31$) and in direction-correct
rate (Fisher $p=1.0$).
Descriptively, Fixed-Warmup tends to over-calibrate (2/3 seeds with
$P(A{<}0){>}0.8$, Fisher $p=0.40$ versus Fixed) while Fixed lands centred
(seed~1 at $P(A{<}0){=}0.744$).
The over-calibration tendency is a candidate hypothesis for follow-up seed
expansion; with $n=3$ per arm, it is not yet established.

\textbf{(8) Fixed seed 1 is the only seed in the study to satisfy all four
classical-quality criteria simultaneously.}
On-target $P(A{<}0)=0.744$, $\mathcal{L}(V)=2.0$ (within $1.8\times$ of
classical), correct $b_\text{final}=378.2$ matching the dataset value scale,
and state-dependent $a_\text{final}=58.5$ within $1.1\times$ of classical's
slope. This is a proof of concept that the DRU architecture can function as
an IQL value network when correctly initialised. The 1/3 (Wilson CI
$[6\%,79\%]$) on-target rate at Fixed (or 1/6 across the Fix-A pool) is
imprecisely estimated; reproducibility requires expansion.

\textbf{(9) A third failure mode identifies the next concrete intervention.}
Fixed seed~0 fails via cold-start value-loss shock:
$\mathcal{L}(V)_0\approx\tau\cdot374^2\approx98\,000$ drives 10\,000
clipped unit-norm gradient steps before $Q$ bootstraps to the correct scale,
causing systematic overshoot ($b_\text{final}=455$).
\textbf{Fix~C}---freezing the value gradient for 1\,000--2\,000 steps to
allow prior $Q$ bootstrapping---is predicted to resolve this without any
architectural change. Fix~C is untested in this report.

\textbf{(10) The diagnostic chain has a constructive resolution, with
provisional statistical support.}
\begin{enumerate}
  \item \emph{Constant-V ablation ($n=3$)}: confirmed $V$ is a passenger in
        the original configuration (83.6\% of classical at zero parameter
        cost).
  \item \emph{No-warmup ablation ($n=8$)}: showed warmup removal alone
        produces returns indistinguishable from Constant-V ($p=0.84$),
        with direction-correct rate 2/8 ($[7\%,59\%]$) and on-target rate
        1/8 ($[2\%,47\%]$).
  \item \emph{Fix-A only (Fixed-Warmup, $n=3$)}: correct scale initialisation
        alone restores competitive return (CV\,=\,1.2\%) and direction-correct
        calibration in 3/3 seeds, with 2/3 over-calibrated.
  \item \emph{Fix-A\,+\,B (Fixed, $n=3$)}: produces seed~1, the only run
        in the study to satisfy all four classical-quality criteria
        simultaneously.
\end{enumerate}
The pooled Fix-A vs default contrast (Welch $p=0.041$, Fisher $p=0.034$) is
the principal statistical evidence.

The remaining failure mode (cold-start overshoot) is mechanistically
understood and has a concrete, untested remedy (Fix~C).
Full statistical validation requires seed expansion to $n\geq8$ for both
Fix-A variants and Fix-C evaluation.
Task-transfer validation to a sparse-reward environment (e.g.\
\texttt{antmaze-medium-play}) where the value function is non-optional
is the essential next step to determine whether these improvements persist
beyond Hopper-medium's forgiving task structure.

%──────────────────────────────────────────────────────────────────────────────
\begin{thebibliography}{9}
\bibitem{kostrikov2021}
I.~Kostrikov, A.~Nair, and S.~Levine,
``Offline Reinforcement Learning with Implicit Q-Learning,''
\textit{ICLR}, 2022.

\bibitem{skolik2021}
A.~Skolik, J.~J.~McClean, M.~Mohseni, P.~van der Smagt, and M.~Leib,
``Layerwise learning for quantum neural networks,''
\textit{Quantum Machine Intelligence}, 3(1):5, 2021.
\end{thebibliography}

\end{document}